% Folder 91, batch 01 — archivists' pages 1–20.
% Pass: Opus 5.5 (claude-opus-5-5), 2026-10-03 — first pass, unchecked against the pages by a human.
%
% Pages 3 and 14 are blank and are absent. Page 1 is a cover in his hand: its
% words become the section heading, with no \page. Page 5 is a first draft of
% the theorem of page 4, cancelled whole and written upside down on the sheet.
\documentclass[11pt,a4paper]{article}
\input{../preamble/grothendieck}

\folder{91}
\batch{1}
\pages{1}{20}
\foldertitle{Autour de Néron / Greenberg-Néron. Foncteurs Hom (méthodes non-projectives) : notes manuscrites (s.d.), lettre (1967).}
\dating{[à partir de 1964-vers 1970]}
\watermark{Édition de démonstration}

\begin{document}

\section{Représentabilité des foncteurs $\underline{\mathrm{Hom}}$, $\prod_{T/S} X/T$, de Greenberg etc.}

\note{Ce titre est le sien : il est seul sur la feuille de la p. 1, qui sert de couverture.}

\page{2}
Représentabilité relative de $\prod_{X/S} P/X$ (sans hyp. projective)

« Flattening functor »

Passage au quotient

$\underline{\mathrm{Hom}}_S(X,Y)$

$\underline{\mathrm{Aut}}_S(X)$

\note{La feuille ne porte que ces lignes, en haut à droite et au milieu : une liste de sujets plutôt qu'un texte.}

\page{4}
\textbf{Th.} Soit $f\colon X\to S$ un morphisme \add{plat} de type fini, avec $S$ loc.\ noeth., $Y\to X$ un morph.\ \add{de type fini} \struck{un sous-préschéma \ill{} de $X$}. Pour que $\prod_{X/S} Y/X$ soit représentable il f.\ et s.\ qu'il commute aux limites projectives adiques d'anneaux artiniens, et \add{si $Y\to X$ est une immersion (fermée)} alors il est représentable par un sous-préschéma (fermé) de $S$.

\marginal{il suffit que $Y\to X$ soit un monomorphisme \ill{}}

\textbf{Dém.} Vérifions les conditions générales du th.\ de représentabilité par morphismes non \uncertain{nécessairement} séparés.

1°) $F$ faisceau pour fpqc : \struck{EGA IV 8} clair.

2°) $F$ commute aux $\varinjlim$ d'anneaux : clair \uncertain{EGA IV 8}.

\marginal{$+$ t.f.}

3°) C'est la condition de commutation aux $\varprojlim$ adiques d'artiniens.

4°), 5°), 6°) Dans le cas d'un monomorphisme $F\to S$ se réduit à ceci : si $A\to A'$ est un hom.\ injectif d'artiniens \add{($A'$ fini sur $A$)}, alors $F(A)\to F(A')$ est bijectif, ou encore $F(A')\neq\emptyset \Rightarrow F(A)\neq\emptyset$. Il faut prouver que si $X_{A'}=Y_{A'}$, alors $X_A=Y_A$. C'est clair, car $X_A$ étant plat sur $A$, $\mathcal{O}_{X_A}\to\mathcal{O}_{X_A}\otimes_A A'$ est injectif.

\marginal{$+$ plat}

7°) Si $T$ sur $S$ est \struck{noeth.} local, noeth., muni d'un idéal de carré nul $J$, $T_0=V(J)$, $U$ ouvert de $T$ tel que $T$ soit l'adhérence schématique de $U$, alors pour \struck{\ill{}} le carré
\[
\begin{array}{ccc}
F(T) & \longrightarrow & F(T_0)\\
\downarrow & & \downarrow\\
F(U) & \longrightarrow & F(U_0)
\end{array}
\]
est cartésien.

\marginal{$+$ plat}

\note{La démonstration de 7° continue en p. 6 ; la p. 5 en est une première rédaction de l'énoncé, biffée.}

\page{5}
\struck{\textbf{Théorème.} Soit $f\colon X\to S$} \add{\struck{plat}}\struck{, de présentation finie, $Y$ un sous-préschéma de $X$ de présentation finie. Pour}

\note{Page entièrement biffée, où l'énoncé s'interrompt après « Pour ».}

\page{6}
\struck{En d'autres termes} \add{r.c.} \ill{} : si $F(T_0)\neq\emptyset$, $F(U)\neq\emptyset$, alors $F(T)\neq\emptyset$. \add{En fait, il suffit ici $F(U)\neq\emptyset$ pour $F(T)\neq\emptyset$} \add{(6 lorsque $Y$ est supposé fermé dans $X$)}. \struck{En d'autres termes si $X_{T_0}=Y_{T_0}$ et} $X_U=Y_U$, alors $X=Y$. En effet, comme $T$ est l'adhérence schématique de $U$ et $X_T$ plat sur $T$, $X_T$ est l'adhérence schématique de $X_U$, donc \struck{on gagne} $Y_T=X_T$ puisque $Y_T$ sous-préschéma fermé de $X_T$ majorant $X_U$.

8°) Si $T$ sur $S$ \uncertain{irréductible} muni de $J$ de carré nul, $T_0=V(J)$, si $F(T_0)\neq\emptyset$, \struck{alors} \struck{$\exists$ ouvert dense} \ill{} $\exists$ ouvert dense $U$ de $T$ \struck{$F(\mathcal{O}_{T,\eta})\neq\emptyset$} \struck{tel que $F$} tel que $F_U$ \uncertain{représentable}.

\marginal{$X$ plat t.f.}

En effet, il suffit que $X_U$ soit « ens.\ libre » sur $U$, et on sait qu'on peut trouver $U$ de cette façon. EGA IV 11.

\add{D}onc \ill{} conditions 1° \ill{} 7° \ill{} sont \uncertain{vérifiées} par le foncteur \ill{}.

\struck{Corollaire}

\struck{\textbf{Corollaire.} Supposons que $X$ soit un $S$-groupe, $Y$ un $S$ sous-groupe. Supposons que $X$ ait la propriété suivante : Pour tout} \struck{$T$} \add{\struck{sur $S$}} \struck{spectre d'un anneau de valuation discrète, et tout} \struck{\ill{} de $X_T$ qui contient \ill{} spécial \ill{}}

\note{Le corollaire est barré de trois longues diagonales ; il est repris p. 7.}

\page{7}
La commutation aux $\varprojlim$ adiques peut aussi s'exprimer ainsi :

(L) Pour tout $T$ \add{sur $S$} spectre d'un anneau local noeth.\ \struck{complet sur $S$} \add{(qu'on peut supposer complet, à corps rés.\ alg.\ clos)}, et tel que $Y_T$ majore un voisinage induit de la fibre spéciale de $X_T$, on a \struck{\ill{}} $X_T=Y_T$.

\textbf{Corollaire.} Cas $X$ propre sur $S$ ; O.K.

Si $X$ est un groupe, $Y$ un sous-groupe, on sait que $Y_T$ \struck{contient} \add{majore} un voisinage induit de la section unité, donc $Y_T\to X_T$ est \add{un mono.} étale, \struck{\ill{}} étant une immersion \struck{fermée}, c'est une immersion ouverte \struck{et fermée}. \add{De plus, la vérification $Y_T=X_T$ se fait} \struck{Donc dans le cas} \uncertain{dans} \struck{\ill{}} « connexité », donc se ramène au cas où $T$ est \add{le spectre d'un anneau} de valuation discrète. \struck{Donc}

\textbf{Corollaire.} Dans ce cas, la condition (L) \uncertain{équivaut} à la suivante : Pour tout $T$ sur $S$ spectre d'un anneau \struck{local noeth.} \add{de valuation discrète} (qu'on peut supposer complet à corps rés.\ alg.\ clos), tel que $Y_T$ soit un sous-groupe ouvert et fermé de $X_T$ qui contient la fibre \struck{générique} \add{spéciale}, on a $Y_T=X_T$. Cette condition est vérifiée dans chacun des cas suivants :

a) $X$ à fibres connexes

b) Il existe une famille de sous-\struck{groupes finis} \add{schémas} $G_i$ de $G$, finis sur $S$, et telle que pour chaque $s\in S$, la réunion ensembliste des $G_{i,\bar s}\cdot G^0_{\bar s}$ soit $G_{\bar s}$.

\note{Dans b) le groupe s'appelle $G$ et non plus $X$ : c'est sur la page.}

\page{8}
\struck{\textbf{Définition.} Un groupe} \add{\struck{(hypothèses)}} \struck{$X$ sur $S$ est dit} \add{\struck{relativement}} \struck{\ill{}}

\textbf{Définition.} On dit qu'une condition \uncertain{R} est satisfaite par un préschéma en groupes $G$ \struck{de type fini} sur $S$ \struck{localement noeth} si pour tout schéma local $T$ sur $S$, et tout sous-groupe ouvert \struck{et fermé} $U$ de $G_T$ qui majore la fibre spéciale, on a $U=G_T$.

\marginal{Condition R}

Lorsque \struck{$G$ est} $G$ est de type fini sur $S$ loc.\ noeth., il suffit de tester avec $T$ spectre d'un anneau de valuation discrète [complet, corps rés.\ alg.\ clos].

\note{La première ligne de la définition est en partie surchargée ; « une condition R est satisfaite par » est une lecture de sens plus que de lettres.}

\textbf{Corollaire.} Soit \struck{$X$} $X$ de \struck{type fini sur} \add{type} \struck{présentation} finie sur $S$ \add{loc.\ noeth}, \struck{$H$} $Y$ un sous-groupe fermé de $X$, \struck{présentation finie}. Supposons que $X$ satisfait R. Alors $\prod_{X/S} Y/X$ est représentable par un sous-préschéma fermé de $S$.

\textbf{Exemples}

\struck{Soit} \struck{\textbf{Corollaire.} Soit $G$ schéma en groupes noeth.\ sur $S$ loc.\ noeth., $H$ un sous-groupe fermé de type fini sur $S$ satisfaisant à la condition R} \add{\struck{(si $G$ séparé sur $S$)}}\struck{. Alors $\mathrm{Cent}_G(H)$ et $\mathrm{Norm}_G(H)$ sont représentables par un sous-préschéma fermé de $G$. De même $\mathrm{Transp}_G(H,K)$.}

\note{Ce dernier corollaire est barré de deux diagonales ; il est repris p. 9.}

\page{9}
\textbf{Corollaire.} Soit $u\colon H\to G$ un homom.\ de préschémas en groupes \add{sur $S$}, \add{et} \struck{séparés sur $S$}, avec $G$ loc.\ noeth., $H$ de type fini sur $S$, \struck{Alors} $H$ \add{plat et} satisfaisant à la condition R \struck{si $G$ séparé sur $S$} (p.\ ex.\ à fibres connexes…). Alors $\mathrm{Cent}_G(u)$ est représentable par un sous-préschéma \struck{fermé} de $G$, \struck{Celui-ci} fermé si $G$ est séparé sur $S$. Si $K$ est un sous-groupe de $G$ \add{\uncertain{de t.f.\ sur $S$}}, \struck{alors} $\mathrm{Transp}_G(u,K)$ est représentable, et si $K$ est un sous-préschéma de $G$ (resp.\ sous-préschéma fermé de $G$) il en est de même de $\mathrm{Transp}_G(u,K)$. Enfin, si \struck{$u$ et $H$} $K$ est \uncertain{également} plat sur $S$ et satisfait la condition R, et si $u$ est un monom., alors \struck{Pour $\mathrm{Transp}_G$} \struck{$\mathrm{Transp}_G(u,K)$} $\mathrm{Transp}_G(H,K)$ est représentable par un \struck{sous-pré}schéma de type fini sur $G$, et c'est un sous-préschéma (resp.\ sous-préschéma fermé) s'il en est de même de $H$ et $K$ dans $G$. En particulier, $\mathrm{Norm}_G(H)$ est représentable, c'est un sous-préschéma (sous-préschéma fermé) de $G$ si $H$ l'est.

\page{10}
\textbf{Variantes} \uncertain{débarrassées} d'hypothèses noethériennes \ill{}

\textbf{Variante.} Pour un morphisme $Y\to X$ non ramifié.

N.B. Il faudrait également \ill{} du Th.\ précédent \ill{}.

\textbf{Théorème.} Soit $f\colon X\to S$ morphisme de type fini, $S$ loc.\ noeth., $E$ cohérent sur $X$. Soit $F$ le sous-foncteur de $S$ :
\[
F(S')=\begin{cases}\emptyset & \text{si } E_{S'} \text{ non plat sur } S'\\ \{\emptyset\} & \text{sinon}\end{cases}
\]
\uncertain{(comme Module plat)}. Pour que ce foncteur soit représentable, il f.\ et s.\ qu'il commute aux limites projectives adiques d'anneaux artiniens.

\uncertain{Vérification en divers pts.}

1°) Faisceau pour fpqc : clair.

2°) Commute aux $\varinjlim$ d'anneaux : EGA IV 8 et IV 11.

\marginal{$X$ type fini sur $S$, $E$ \uncertain{coh.}}

3°) Condition de commutation aux $\varprojlim$ d'artiniens.

4°) 5°) 6°) Si $A\to A'$ injectif ($A'$ fini sur $A$ local artinien) alors $F(A)\to F(A')$ bijectif, i.e.\ $F(A')\neq\emptyset\Rightarrow F(A)\neq\emptyset$, i.e.\ $X_{A'}$ plat \uncertain{(je veux dire $E$ plat)} sur $A'$ $\Rightarrow$ $X_A$ plat \uncertain{(id.)} sur $A$. EGA IV 11.

7°) \ill{} $F(T_0)\neq\emptyset$, $F(U)\neq\emptyset \Rightarrow F(T)\neq\emptyset$, i.e.\ $X_{T_0}$ plat et $X_U$ plat $\Rightarrow$ $X_T$ plat. EGA IV 11.

\note{Ici comme en 4°–6°, il écrit $X_\cdot$ plat là où il s'agit de la platitude de $E_\cdot$ ; il le dit lui-même en 4°.}

\page{11}
8°) \ill{} Trouver $U$ ; il suffit de trouver que $X_{U_0}$ soit « ens.\ libre » sur $U_0$, OK par EGA IV 11.

La condition énoncée dans le th.\ signifie ceci :

(L) Pour tout $T$ spectre d'un anneau \add{local noeth.} \struck{de valuation discrète}, $\{$complet à corps résiduel alg.\ clos si on y tient$\}$, tel que $E_T$ sur $X_T$ soit plat sur $T$ dans un voisinage de la fibre \uncertain{spéciale} \ill{}, \struck{alors} $E_T$ est plat sur $T$. \struck{Exemple}

\marginal{N.B. On peut supposer $\dim T=1$, \ill{} $T$ irréductible \ill{} i.e.\ \ill{} spectre d'un anneau de valuation discrète \ill{}}

\textbf{Corollaire 1.} O.K.\ si $f\colon X\to S$ est propre.

\textbf{Corollaire 2.} O.K.\ si \struck{$G$} $E=\mathcal{O}_X$, $X$ un préschéma en groupes sur $S$, satisfaisant à la condition R plus haut.

\textbf{Corollaire 3.} Soit $f\colon X\to S$ un morphisme de \struck{type fini} \add{prés.\ finie}, $S$ \struck{loc.\ noeth.} \ill{}. Conditions équivalentes \struck{\ill{}} :

(i) $f$ se factorise en $f'f''$, $X\xrightarrow{f''}S'\xrightarrow{f'}S$, avec $f''$ fid.\ plat \uncertain{et quasi-compact}, $f'$ immersion

(i bis) id., avec $f''$ plat, $f'$ imm.,

(ii) \uncertain{Le morphisme} $X\times_S X\xrightarrow{\mathrm{pr}_1}X$ est plat.

\note{La disposition de la factorisation $X\to S'\to S$ en (i) est restituée ; le second mot après « fid.\ plat » est peu lisible.}

\page{12}
Sous ces conditions, \struck{la} factorisation (i) est unique à isom.\ près, et $S'$ est de \struck{t.f.} \add{prés.\ finie} sur $S$.

(i) $\Rightarrow$ (i bis) $\Rightarrow$ (ii) trivial. D'où l'unicité de la factorisation (i), car $X\to S'$ sera quotient, donc quotient pour fpqc, donc $S'=X/R$, où $R=X\times_S X$.

Reste à prouver (ii) $\Rightarrow$ \add{le faisceau} $X/R$ est représentable par un $S'$ de \add{prés.} \struck{type} finie sur $S$. \struck{(N.B. alors, \ill{})}

[N.B. Alors $X\to S'$ est couvrant pour fpqc, comme on le voit \ill{} $X\to S'$ et donc $X\times_S X=R\to X$ qui est fid.\ plat par hypothèse, il est fid.\ plat \ill{}, ce qui implique \add{aussi} que $S'\to S$ est \add{\ill{}} \struck{donc plat}.]

Or \struck{\ill{}} le foncteur représenté par $X/R\to S$ est \add{(i.e.\ le suivant :} \add{$\Leftrightarrow X_T\to T$ couvrant fpqc)}
\[
F(T)=\begin{cases}\emptyset & \text{si } X_T \text{ non fid.\ plat sur } T\\ \{\emptyset\} & \text{sinon}\end{cases}
\]
D'ailleurs on est ramené au cas $S$ affine, noethérien. \struck{et} Alors on applique le th., qui nous dit qu'il suffit de vérifier \struck{pour} (L). Or \ill{} de cette condition, $X_T\to T$ est couvrant [car il existe un pt $x$ de la fibre spéciale, et pour \ill{} $\mathcal{O}_{X,x}$ sur $\mathcal{O}_{T,t_0}$ est fid.\ plat] \ill{}.

\page{13}
\textbf{Exemple}

\textbf{Corollaire.} Soit $G$ un préschéma en groupes de \struck{type} \add{prés.} finie sur $S$, \struck{loc.\ noeth.} $P$ un \struck{fibré principal} préschéma de prés.\ finie sur $S$ sur lequel $G$ opère, de façon que $P$ soit formellement principal homogène sous $G$. Sous ces conditions, il existe un sous-préschéma $S'\hookrightarrow S$, tel que $P\to S$ se factorise par $P\to S'$, et que \struck{$P$} $P$ devienne \struck{formellement} principal homogène sous $G$, i.e.\ $P\to S'$ est \struck{plat} fid.\ plat. [De plus, $S'$ est unique à isomorphisme unique près, et est de prés.\ finie sur $S$.]

\textbf{Corollaire.} Soit $u\colon G\to H$ un homom.\ de $S$-préschémas en groupes, \struck{de présentation} et de présentation finie. Supposons $\mathrm{Ker}\,u=N$ [qui est de prés.\ finie sur $S$] plat sur $S$. Alors \struck{le foncteur} $G/N$ existe.

\page{15}
\begin{tikzcd}[column sep=small, row sep=small]
Z \arrow[r, "\varphi"] \arrow[dr, "h"'] & X \arrow[d, no head, "f"] & Y \arrow[dl, no head, "g"] \\
 & S &
\end{tikzcd}

\begin{tikzcd}[column sep=small, row sep=small]
{Z'} \arrow[r, "{\varphi'}"] \arrow[rr, bend left=20, "v"] \arrow[dr, "{h'}"'] & {X'} \arrow[r, "w"] \arrow[d, "{f'}"] & {Y'} \arrow[dl, "{g'}"] \\
 & {S'} &
\end{tikzcd}

\note{Dans le premier diagramme, $f$ et $g$ sont deux traits sans pointe ; dans le second, toutes les flèches vont vers $S'$. L'étiquette de la flèche $X'\to Y'$ se lit $w$, celle de l'arc $Z'\to Y'$ se lit $v$.}

\marginal{Énoncé plutôt en termes de $\prod_{X/S} P/X\to\prod_{Z/S} Q/Z$}

$f$, $h$ plat et propre sur $S$ loc.\ noeth.

$g$ \struck{$g,h$} de type fini [et \uncertain{non ramifié} ?]

\marginal{suffit si $Z\to X$ est surjectif}

N.B. L'hypothèse $S$ loc.\ noeth.\ peut être remplacée par des hypothèses de présentation finie sur $f,g,h$.

Posons \ill{}

a) Le sous-foncteur de $\underline{\mathrm{Hom}}_S(X,Y)$ \ill{} $u\mapsto u\circ\varphi$, $\underline{\mathrm{Hom}}_S(X,Y)\to\underline{\mathrm{Hom}}_S(Z,Y)$ \uncertain{non ramifié} sur \add{$\underline{\mathrm{Hom}}_S(Z,Y)$} \ill{}, est un sous-foncteur \uncertain{ouvert}.

Il correspond à $u'\colon X'\to Y'$ tel que, pour tout $s\in S'$,
\[
\mathrm{Hom}\bigl(u_s'^{*}(\Omega^1_{Y'_s/k(s)}),k(s)\bigr)\longrightarrow\mathrm{Hom}\bigl(v_s^{*}(\Omega^1_{Y'_s/k(s)}),k(s)\bigr)
\]
soit injectif. Or on peut trouver sur $S'$ des Modules \uncertain{cohérents} $M_u$, $M_v$ tels que pour \uncertain{tout} $J$ \uncertain{quasi-cohérent} sur $S'$, on ait des isomorphismes \uncertain{fonctoriels}
\[
\mathrm{Hom}_{\mathcal{O}_{X'}}\bigl(u^*(\Omega^1_{Y'/S'}),\mathcal{O}_{X'}\otimes_{\mathcal{O}_{S'}}J\bigr)\simeq\mathrm{Hom}_{\mathcal{O}_{S'}}(M_u,J),
\]
\[
\mathrm{Hom}_{\mathcal{O}_{Z'}}\bigl(v^*(\Omega^1_{Y'/S'}),\mathcal{O}_{Z'}\otimes_{\mathcal{O}_{S'}}J\bigr)\simeq\mathrm{Hom}_{\mathcal{O}_{S'}}(M_v,J),
\]
et la formation de $M_u$, $M_v$ \uncertain{commute} \ill{} à l'extension de la base. Ceci dit, \ill{} un homomorphisme
\[
(*)\qquad M_v\longrightarrow M_u
\]
qui définit par $\mathrm{Hom}(\;,J)$ l'homomorphisme naturel
\[
\mathrm{Hom}_{\mathcal{O}_{X'}}\bigl(u^*(\Omega^1_{Y'/S'}),\mathcal{O}_{X'}\otimes_{\mathcal{O}_{S'}}J\bigr)\longrightarrow\mathrm{Hom}_{\mathcal{O}_{Z'}}\bigl(v^*(\Omega^1_{Y'/S'}),\mathcal{O}_{Z'}\otimes_{\mathcal{O}_{S'}}J\bigr).
\]
Dire \struck{\ill{}} la condition de non \uncertain{ramification} \ill{} en un pt $s\in S'$, signifie que
\[
M_v\otimes k(s)\longrightarrow M_u\otimes k(s)
\]
\uncertain{est surjectif}, et par Nakayama ceci \ill{}

\marginal{N.B. \ill{}}

\note{Page très serrée en haut ; les deux diagrammes, l'encadré et les hypothèses y sont disposés côte à côte, et sont transcrits l'un après l'autre. La phrase continue en tête de la p. 16.}

\page{16}
voisinage.

Nous désignerons par
\[
F=F_\varphi=\underline{\mathrm{Hom}}_S(X,Y;\varphi)
\]
le sous-foncteur précédent de $\underline{\mathrm{Hom}}_S(X,Y)$, dont les valeurs sur $S'$ sont l'ens.\ des $S'$-homomorphismes $u\colon X'\to Y'$ tels que, posant $v=u\varphi'$, et considérant $M_v$, $M_u$, $M_v\to M_u$ correspondants, $M_v\to M_u$ soit surjectif.

\textbf{Théorème.} $\underline{\mathrm{Hom}}_S(X,Y;\varphi)\to\underline{\mathrm{Hom}}_S(Z,Y)$ est représentable \uncertain{par immersions ouvertes} \ill{}.

\textbf{Corollaire.} Si $\underline{\mathrm{Hom}}_S(Z,Y)$ est représentable, alors $\underline{\mathrm{Hom}}_S(X,Y;\varphi)$ l'est également, et le morphisme $\underline{\mathrm{Hom}}_S(X,Y;\varphi)\to\underline{\mathrm{Hom}}_S(Z,Y)$ est \uncertain{non ramifié}.

\textbf{Principe de la démonstration.} On peut supposer donné un $v\colon\underline{\mathrm{Hom}}_S(Z,Y)$ \ill{}, et on doit représenter le sous-foncteur de $\underline{\mathrm{Hom}}_S(X,Y)$ dont les valeurs sur $S'$ sont l'ens.\ des homomorphismes $u'\colon X'\to Y'$ tels que 1°) $u'\circ\varphi'=v'$ 2°) La condition $M_{v'}\to M_{u'}$ surjectif. On applique le critère général des représentabilités \add{\ill{}} \struck{foncteurs}

\marginal{N.B. Réduction au cas $S$ \uncertain{noeth.\ intègre} \ill{} dans toute \ill{} représentabilité \ill{}}

\page{17}
préschémas \struck{\ill{}} \add{\uncertain{non ramifiés}} \struck{\ill{}} \uncertain{séparés} sur $S$. Seules \add{\uncertain{deux conditions}} les conditions de \uncertain{modularité} ne résultent pas des techniques habituelles. \struck{\ill{}} Celles qui \ill{} se \uncertain{démontrent} par \ill{}, avec le calcul habituel d'obstructions.

\marginal{\ill{}}

Reste le cas où $S$ est local \ill{} de dim 1, avec Idéal $J$ annulé par \uncertain{l'idéal maximal}, d'où $S_0=V(J)$, et où on dispose de $u_0\colon X_0\to Y_0$ (tel que $u_0\varphi_0=v_0$) et un prolongement $u_\eta\colon X_\eta\to Y_\eta$ ($\eta$ le pt générique de $S$) tel que $u_\eta\varphi_\eta=v_\eta$, $u_{\eta 0}=u_{0\eta}$. Question de prolonger $u_0$ en $u$ tel que $u\varphi=v$. Soient $P(u_0,S)$, $P(u_0,S_\eta)$ ; $P(v_0,S)$, $P(v_0,S_\eta)$ les prolongements de $u_0$ resp.\ $v_0$ sur $S$, \struck{\ill{}} de $u_{0\eta}$ resp.\ $v_{0\eta}$ sur $S_\eta$, d'où le diagramme \uncertain{commutatif}

\begin{tikzcd}[column sep=small, row sep=small]
\text{(1)} & P(u_0,S) \arrow[r] \arrow[d] & P(u_0,S_\eta)\ast \arrow[d] \\
 & \ast P(v_0,S) \arrow[r] & P(v_0,S_\eta)
\end{tikzcd}

d'ensembles principaux homogènes sous les gr.\ du diagramme

\begin{tikzcd}[column sep=small, row sep=small]
\text{(2)} & \mathrm{Hom}(M_{u_0}^{*},J) \arrow[r] \arrow[d] & \mathrm{Hom}(M_{u_0},J_\eta) \arrow[d] \\
 & \mathrm{Hom}(M_{v_0},J) \arrow[r] & \mathrm{Hom}(M_{v_0},J_\eta)
\end{tikzcd}

Par hypothèse on tient des éléments aux \ill{}

\note{L'indice qu'il écrit comme un $y$ est lu $\eta$, le point générique qu'il nomme ainsi plus haut. Les deux étoiles du diagramme (1) marquent les éléments dont parle la p. 18 ; l'astérisque de $M_{u_0}^{*}$ dans (2) est sur la page.}

\page{18}
marqués de $(\ast)$, dont l'image dans le coin \uncertain{droit du bas} sont les mêmes, on cherche un élément du coin \uncertain{gauche} \uncertain{en haut} qui \ill{} nos éléments marqués. On note qu'il suffit pour ceci que $P(u_0,S)\neq\emptyset$, car alors le diagramme (1) est isomorphe à (2), or (2) est \emph{cartésien}, grâce au fait que $M_{v_0}\to M_{u_0}$ est surjectif. Or l'obstruction à trouver un élément dans $P(u_0,S)$ est \struck{\ill{}}
\[
\partial(u_0)=\xi\in\mathrm{Ext}^1_{\mathcal{O}_{P_0}}(P_0;I_{\Gamma_0},\mathcal{O}_{\Gamma_0}\otimes_{S_0}J)
\]
où $P_0=X_0\times_{S_0}Y_0$, \struck{\ill{}} $\Gamma_0$ est le sous-préschéma de $P_0$ graphe de $u_0$, et où $I_{\Gamma_0}$ est l'Idéal de $P_0$ qui le définit. C'est un module sur $A_0$, qui s'annule quand on \struck{\ill{}} \add{localise en $\eta$}, donc si $\pi$ est une \uncertain{uniformisante} de $A$, il existe un $n\geqslant 0$ tel que $\pi^n\xi=0$. Utilisons la suite exacte
\[
(\ast)\qquad 0\to J\xrightarrow{\;\pi^{n}\;}J\to J_n\to 0
\]
et la suite \uncertain{exacte} correspondante
\[
0\to\mathcal{O}_{\Gamma_0}\otimes_{S_0}J\xrightarrow{\;\pi^{n}\;}\mathcal{O}_{\Gamma_0}\otimes_{S_0}J\to\mathcal{O}_{\Gamma_0}\otimes_{S_0}J_n\to 0
\]
et la suite exacte des $\mathrm{Ext}^i_{\mathcal{O}_{P_0}}(P_0;I_{\Gamma_0},-)$ ; on trouve que $\xi$ provient d'un élément
\[
\xi'\in\mathrm{Hom}_{\mathcal{O}_{P_0}}(P_0;I_{\Gamma_0},\mathcal{O}_{\Gamma_0}\otimes_{S_0}J_n)
\]
\uncertain{mod} $\mathrm{Im}\,\mathrm{Hom}_{\mathcal{O}_{P_0}}(P_0;I_{\Gamma_0},\mathcal{O}_{\Gamma_0}\otimes_{S_0}J)$.

\marginal{\ill{} $P_0$ \ill{} $P_1$ \ill{}}

\note{L'exposant de $\pi$ au-dessus de la flèche de $(\ast)$ est à demi effacé ; $\pi^{n}$ est la lecture qui s'accorde avec $\pi^n\xi=0$.}

\page{19}
Or le \add{premier terme des} \struck{deuxième} membres se réduit à
\[
\mathrm{Hom}_{\mathcal{O}_{P_0}}(I_{\Gamma_0}/I_{\Gamma_0}^2,\mathcal{O}_{\Gamma_0}\otimes_{S_0}J_n)\simeq\mathrm{Hom}_{\mathcal{O}_{X_0}}\bigl(u_0^*(\Omega^1_{Y_0/S_0}),\mathcal{O}_{X_0}\otimes_{S_0}J_n\bigr)
\]
\[
\simeq\mathrm{Hom}_{S_0}(M_{u_0},J_n)
\]
et de même \struck{\ill{}} le deuxième terme s'écrit $\mathrm{Hom}_{S_0}(M_{u_0},J)$, de sorte que l'obstruction peut être considérée comme un élément
\[
\xi''\in\mathrm{Hom}_{A_0}(M_{u_0},J_n)/\mathrm{Im}\,\mathrm{Hom}_{A_0}(M_{u_0},J)\hookrightarrow\mathrm{Ext}^1_{A_0}(M_{u_0},J)
\]
[en utilisant encore, pour la dernière inclusion, la suite exacte $(\ast)$]. Donc \uncertain{on obtient} un élément
\[
\xi''\in\mathrm{Ext}^1_{A_0}(M_{u_0},J)
\]
annulé par une puissance de $\pi$ dont l'annulation est nécessaire et suffisante pour la possibilité de prolonger $u_0$ en $u\colon X\to Y$. On trouve de même un
\[
\eta''\in\mathrm{Ext}^1_{A_0}(M_{v_0},J)
\]
et nous admettrons que $\eta''$ est l'image de $\xi''$ par $M_{v_0}\to M_{u_0}$. Or \emph{$\eta''$ est nul}, \uncertain{car} $v_0$ se prolonge en $v$, \ill{}. Utilisons la suite exacte
\[
0\to N\to M_{v_0}\to M_{u_0}\to 0,
\]
il vient, par la suite exacte des $\mathrm{Ext}^i(\;,J)$, que $\xi''$ correspond à un élément

\note{« $\eta''$ est nul » est souligné par lui.}

\page{20}
\[
\xi'''\in\mathrm{Hom}(N,J)/\mathrm{Im}\,\mathrm{Hom}(M_{v_0},J).
\]
\struck{Cet} Cet élément localisé en $\eta$ est nul, \struck{\ill{}} mais c'est là tout ce que nous pouvons tirer du fait que $u_0\varphi_0$ et $u_\eta$ se prolongent. Mais \add{$v$ et $u_\eta$} nous savons plus : les prolongements \uncertain{sont faits} de façon que $v_\eta=$ \struck{$\varphi_\eta$} $u_\eta\varphi_\eta$. À priori, l'idée une fois acquise \uncertain{le jeu} des 2 possibilités \emph{séparées} des deux obstructions, i.e.\ dans (1) de l'existence de deux éléments dans les \uncertain{coins}, la possibilité de les choisir de façon \uncertain{cohérente} se \uncertain{trouve} dans
\[
\mathrm{Hom}(M_{v_0},J_\eta)/\mathrm{Im}\bigl(\mathrm{Hom}(M_{v_0},J)+\mathrm{Hom}(M_{u_0},J_\eta)\bigr)
\]
\[
=\mathrm{Hom}(N,J_\eta)/\mathrm{Im}\,\mathrm{Hom}(M_{v_0},J).
\]
Soit $\eta'$ cet élément. Notons que $\xi'''$ définit également un élément de ce groupe, puisque
\[
\mathrm{Hom}(N,J_n)/\mathrm{Im}\,\mathrm{Hom}(M_{v_0},J)\subset\mathrm{Hom}(N,J_\eta)/\mathrm{Im}\ldots
\]
Nous admettrons que l'on a $\eta'=\xi'''$. Comme par hypothèse $\eta'=0$, on a $\xi'''=0$, OK.

\note{Dans ces deux pages les indices $J_n$ et $J_\eta$ se ressemblent sur la feuille ; la lecture suit la p. 18, où $J_n$ est le quotient de $J$ par $\pi^{n}$, et $J_\eta$ est le localisé en $\eta$.}

\end{document}
